\section{Strict weighted transport}\label{sec:transport}

The pairing will imply \(\E J_F(i)\le0\), with equality only if
\(J_F(i)=0\) almost surely. Comparing this conclusion with
Equation~\eqref{eq:Jacobian} will reduce the contradiction to a
rigidity argument.

Put \(J_{F,\pm}=\max(\pm J_F,0)\) and
\(\omega_F(z)=y^{\beta-1}|Q_F(z)|\).
For any fixed \(F\) and nonnegative Borel \(h\), local change of variables
gives
\begin{equation}\label{eq:area-formula}
 \int_\HH h(z)\omega_F(z)J_{F,\pm}(z)\,dA(z)
 =
 \int_\C\!
 \sum_{\substack{G_F(z)=\xi\\ \pm J_F(z)>0}}
 h(z)y^{\beta-1}|Q_F(z)|^3\,dA(\xi).
\end{equation}
Indeed cover the indicated sign locus by countably many local
diffeomorphism charts, disjointify into Borel pieces, and use absolute
Jacobian \(|J_F||F'|^2\). No global properness of \(G_F\) is required.
At a preimage, \eqref{eq:critical-value} and \(3(k-1)=\beta-1\) give
\begin{equation}\label{eq:target-weight}
 y^{\beta-1}|Q_F(z)|^3=(V_\xi^F(z)/k)^3.
\end{equation}

\begin{lemma}[Goodness at the normalized point]\label{lem:base-good}
The positive-Jacobian mass at the base point is supported on good fibers:
\begin{equation}\label{eq:base-good}
 \E\bigl[J_{F,+}(i)
       \ind_{\{(F,G_F(i))\text{ not good}\}}\bigr]=0.
\end{equation}
\end{lemma}
\begin{proof}
For almost every \(F\), bad targets have area zero by
Lemma~\ref{lem:good}. Formula~\eqref{eq:area-formula} therefore gives
\[
 \int_{\{|x|\le1,\ 1\le y\le2\}}
 \omega_F(z)J_{F,+}(z)
 \ind_{\{(F,G_F(z))\text{ not good}\}}\,dA(z)=0.
\]
The product of \(J_{F,+}\) and the indicator equals its base-point
version for \(T_zF\), by \eqref{eq:G-action}--\eqref{eq:V-action}.
Take expectations and apply \eqref{eq:law} and Tonelli.
The spatial density is \(dx\,dy/y\); its integral on this box is
\(2\log2>0\). This proves \eqref{eq:base-good}.
\end{proof}

The weighted assertion is essential: \(G_F(i)=-i/k\) for every \(F\),
so area-genericity alone says nothing about this fixed target.
We need only the weighted conclusion, not almost-sure goodness
at the base point without the Jacobian factor.

\Needspace{10\baselineskip}
\begin{proposition}[Strict weighted transport]\label{prop:strict-transport}
Assume \(\beta\ge1/4\), let \(P\) be the law of
Proposition~\ref{prop:law}, and set \(k=(\beta+2)/3\).
Then \(\E J_F(i)\le0\).
Equality is possible only if \(J_F(i)=0\) almost surely.
\end{proposition}
\begin{proof}
For \(N\ge1\), restrict the pairing domain \(\mathcal D\) to
\(\mathcal D_N\) by requiring, for \(z'=R_F(z)=x'+iy'\),
\[
 1/N\le y'/y\le N,\qquad |x'-x|\le Ny.
\]
These Borel sets increase to \(\mathcal D\), and their indicators
at \((F,z)\) equal their indicators at \((T_zF,i)\).
Define the equivariant Borel ratio
\[
 c_F(z)=
 \left(\frac{V_{G_F(z)}^F(R_F(z))}{V_{G_F(z)}^F(z)}\right)^3>1
 \quad\text{on }\mathcal D,
\]
and set it to zero elsewhere.
For \(S>0\), \(Y>1\), use the boxes
\[
 B=\{|x|\le S,\ 1\le y\le Y\},\qquad
 B'=\{|x|\le S+NY,\ 1/N\le y\le NY\}.
\]
Every included saddle in \(B\) has its partner in \(B'\).
On the target side of \eqref{eq:area-formula}, multiplication
by \(c_F\) makes its weight exactly the partner's weight
by \eqref{eq:target-weight}. Fiberwise injectivity of the pairing
therefore proves, for every \(F\),
\[
 \int_B\omega_FJ_{F,+}c_F\ind_{\mathcal D_N}(F,z)\,dA(z)
 \le\int_{B'}\omega_FJ_{F,-}\,dA(z).
\]
Take expectations and apply \eqref{eq:law} to the equivariant
nonnegative tests. The resulting inequality is
\begin{align*}
 2S\log Y\,
 \E[J_{F,+}(i)c_F(i)\ind_{\mathcal D_N}(F,i)]
 \le
 2(S+NY)(\log Y+2\log N)\,\E J_{F,-}(i).
\end{align*}
Here \(J_F(i)\) is bounded by compact-family continuity, so
the right side is finite. At fixed \(N\), first let \(S\to\infty\),
then \(Y\to\infty\). Finally let \(N\to\infty\) by monotone
convergence. This gives
\[
 A:=\E[J_{F,+}(i)c_F(i)\ind_{\mathcal D}(F,i)]
 \le d:=\E J_{F,-}(i).
\]
Lemma~\ref{lem:base-good} gives
\(b:=\E J_{F,+}(i)=\E[J_{F,+}(i)\ind_{\mathcal D}(F,i)]\).
Thus \(b\le A\le d\), proving \(\E J_F(i)\le0\).
If equality holds, then \(b=d\) and
\[
 0\le\E[J_{F,+}(i)(c_F(i)-1)\ind_{\mathcal D}(F,i)]
 =A-b\le0.
\]
The integrand is strictly positive wherever \(J_{F,+}(i)>0\)
on \(\mathcal D\). Equation~\eqref{eq:base-good} removes positive
mass outside this domain, so \(b=0\), and hence \(d=0\).
No uniform positive gap in \(c_F-1\) is needed.
\end{proof}
